%\addcontentsline{toc}{part}{Introduction}
\chapter{Introduction}

%	%\addcontentsline{toc}{chapter}{Motivation}
\chapter{Motivation}

Robotics is a field that is experiencing a huge impulse currently. And more specifically, Open Source robotics is growing rapidly. 
\\

The idea of creating common software, of aiding other investigators to easily replicate the work already done and continue from that point instead of "reinventing the wheel" started in the past century. In 1998 the Open Source Initiative (OSI)\cite{osi} was formed, and its definition is recognized as the standard. According to them, "Open source does not just mean access to the source code. The distribution terms of open-source software must comply with the following criteria: Free Redistribution, Source Code, Derived Works, Integrity of The Author's Source Code, No Discrimination Against Persons or Groups, No Discrimination Against Fields of Endeavor, Distribution of License, License Must Not Be Specific to a Product, License Must Not Restrict Other Software,  License Must Be Technology-Neutral"\cite{osi_def}. 
\\

It is my believe that Open Source is critical in the development of new knowledge and not only in the Software field, but in all the science and technical disciplines. This is how the idea of increasing that common pool of tools with another one that might be useful appeared. 
\\
\begin{wrapfigure}{r}{0.3\textwidth}
	\centering
    \includegraphics[width=0.25\textwidth]{img/intro/open_source.eps}
	\caption[Open Source Initiative Logo]{Open Source Initiative Logo}
\end{wrapfigure}

It is noticeable the mark that Open Source has throughout the project. From the libraries being used in it to the Robotic Operating System of which this project is but another package. All of them are Open Source. Also, the tutorials, examples and different web-pages with useful comments and aids for those who are learning has been critical in this project's development. Nothing of this could have been possible without the idea of Open Source code. 
\\

Apart from the want of making useful Open Source code, there remain another motivation unexplained. That is, the specific subject of this project. 
\\

Previously, I have developed other projects involving designing and constructing wheeled robots. It was then the idea of creating different packages that glued together could suit almost every robot appeared. One of the most interesting fields that can be used in virtually all types of robots is Computer Vision. 
\\

Computer vision has experience an important improvement in the last years through the upgrade of the hardware and software that compose it. The hardware such as acquisition elements (cameras, depth sensors, etc) and computing elements (PCs or other programmable devices) have experienced a rapid advance in the past years. It allowed to process more data that is now obtained more accurately and with less noise. This increase in the computing power of the equipment created a possibility of introducing more complex libraries and frameworks and even operating systems. 
\\

But Computer Vision is an ample area, why choosing object recognition? The project being developed is intended to be useful, not just a mere hobby. It occurred to me that apart from being useful in the interaction with robots and cognitive robots more particularly, this software could be used with visually impaired people, as an example. Mainly, with people that had a sudden loss of vision and that is not used to detect the objects by its shape, or cannot differentiate them by the contour. 

 
%	%\addcontentsline{toc}{chapter}{Socio-economical context}
\section{Socio-economical context}
\label{context}
Europe's population is ageing.
This phenomena occurs because of the augment in the live expectancy and the declining of birth rates. 
Recent studies provided by the World Health Organization (WHO) showed that the live expectancy in Spain is the highest in Europe, specially in women \cite{who_live_expectancy}.
Spanish women live an average of  85.1 years in 2012, and the men around 79 years. 
In fact, the Spanish women live expectancy is only behind the Japanese women, that is, the second highest level in the world. 
\\

In the World Population Ageing Report of 2013 it is stated that "by 2050, Bosnia and
Herzegovina, Germany, Malta, Portugal, Serbia and Spain are projected to attain median ages of
50 years or more"\cite{world_population_ageing}. 
This means that soon enough there are going to be more elder people than young people. 
How are we going to be able to support and aid that major sector of the population?
The solution might be to incorporate robots to help in this task. 
\\

Assistive robots can help elder people to interact with their environment. 
They can cooperate in the health control of the user testing the blood pressure or recognizing falls. 
Also, social robots may maintain a conversation and interact with people, for example, reminding them to take a specific medication. 
Japan is the country with the highest population ageing of the world. 
This country is the world leader in the assistive robotics field. 
Recently the he International Organization for Standardization (ISO) has adopted Japanese standards for assistive robotics technology \cite{japan1}\cite{japan2}.
The Japanese government is studying to introduce a system of assistive technologies and hence the funding in this field is high. 
As an example of the technology being developed in this country, the Pepper assistive robot was recently anounced by SoftBank Corp\cite{cheap_humanoid}.
It is a robot that costs less than two thousand USD and that is capable of maintaining a conversation with the user. 
Pepper is not capable nowadays of performing assistive tasks. 
It has a humanoid form with arms, but they are only used to express and empathize the information that it is giving. 
Is hence the first step towards more environment-aware and interactive low-cost robotics.
% Is humanoid and has arms, but at least until now they are only used to express and empathize its words.  
\\

In Spain the assistive field is mainly being evolved in universities and investigation groups. 
Projects such as the RobAlz between the Fundación Alzheimer España (FAE) and the Carlos III University \cite{robalz} are developed to improve the interaction between robots and, in this particular case, Alzheimer patients.  
Another example of the Spanish investigation in the field is the work performed by the Ave María Foundation \cite{assistive_spain}.
They have introduced various robots in the therapies of their patients. 
Nao is a small robot that can recognize faces, sounds and is able to walk. 
It is used to instruct therapy seasons. 
Another one is the robot REEM, which is more than 1.6 m height. 
It is used to transport the necessary items such as clothes or food around the place, allowing the employees to center in assistive tasks. 



% In the past seven years the world has experienced the effects of the recession that appeared in the late 2000s decade. It still affects the economy nowadays and the experts do not know how many more years the situation will stay the same. \\

% In Spain the recession has had a big impact that can still be seen in the unemployment numbers and jobs offers. Also, the poverty has increased enormously and the quality of life has decreased together with the salaries.  
% \\

% There is a general discontent with the politics and the banks. Pessimism is everywhere, as well as the feeling that it is needed a change of direction in both fields. 
% \\


% This economic context has created and motivated new ways of investigating and living, always aiming at spending as less as possible. It has triggered a change in the perception of the knowledge and whether it should be restricted using patents or not. Most of the scientific community now supports the open source initiative, a fact that has increased the rapidity of the new discoveries and technology improvements. 

% \begin{figure}[h]
% 	\begin{center}
%     \includegraphics[scale=0.3]{img/intro/osi.eps}
% 	\caption[Open Source Initiative Logo]{Open Source Initiative Logo}
% 	\end{center}
% \end{figure}


% The idea of creating common software, of aiding other investigators to easily replicate the work already done and continue from that point instead of "reinventing the wheel" started in the past century. In 1998 the Open Source Initiative (OSI)\cite{osi} was formed, and its definition is recognized as the standard. According to them, "Open source does not just mean access to the source code. The distribution terms of open-source software must comply with the following criteria: Free Redistribution, Source Code, Derived Works, Integrity of The Author's Source Code, No Discrimination Against Persons or Groups, No Discrimination Against Fields of Endeavor, Distribution of License, License Must Not Be Specific to a Product, License Must Not Restrict Other Software,  License Must Be Technology-Neutral"\cite{osi_def}. 
% \\

% Open Source is crucial in the development of new knowledge and not only in the Software field, but in all the science and technical disciplines. 




%	%\addcontentsline{toc}{chapter}{Future applications}
\chapter{Future applications}
	% future applications



%\section{Introduction}
%\vspace*{8cm}


\section{Context and motivation}
In the past fifteen years, there has been a change in the perception of the knowledge and whether it should be restricted using patents or not. Most of the scientific community now supports the open source initiative and this has caused a huge increase on the investigation field. \\



The idea of creating common software, of aiding other investigators to easily replicate the work already done and continue from that point instead of "reinventing the wheel" started in the past century. In 1998 the Open Source Initiative (OSI)\cite{osi} was formed, and its definition is recognized as the standard. According to them, "Open source does not just mean access to the source code. The distribution terms of open-source software must comply with the following criteria: Free Redistribution, Source Code, Derived Works, Integrity of The Author's Source Code, No Discrimination Against Persons or Groups, No Discrimination Against Fields of Endeavor, Distribution of License, License Must Not Be Specific to a Product, License Must Not Restrict Other Software,  License Must Be Technology-Neutral"\cite{osi_def}. 
\\

It is my believe that Open Source is critical in the development of new knowledge and not only in the Software field, but in all the science and technical disciplines. This is how the idea of increasing that common pool of tools with another one that might be useful appeared. 
\\
%\begin{wrapfigure}{r}{0.3\textwidth}
%	\centering
%   \includegraphics[width=0.25\textwidth]{img/intro/open_source.eps}
%	\caption[Open Source Initiative Logo]{Open Source Initiative Logo}
%\end{wrapfigure}

\begin{figure}[h]
	\begin{center}
    \includegraphics[scale=0.2]{img/intro/open_source.eps}
	\caption[Open Source Initiative Logo]{Open Source Initiative Logo}
	\end{center}
\end{figure}


It is noticeable the mark that Open Source has throughout the project. From the libraries being used in it to the Robotic Operating System of which this project is but another package. All of them are Open Source. Also, the tutorials, examples and different web-pages with useful comments and aids for those who are learning has been critical in this project's development. Nothing of this could have been possible without the idea of Open Source code. 
\\

%Apart from the want of making useful Open Source code, there remains another motivation unexplained. That is, the specific subject of this project. 
\\

%The bad economic context that exist nowadays has forced the different areas of the society to investigate new technologies that has a lower cost. 
%Nowadays, advances regarding robotics appear everyday, and technologies such as 3D printing are having a tremendous impact in the society. Thanks to this instruments, now it is possible and relatively easy to customize prosthesis and also much affordable than it was before. 
%\\

%Robotics is a field that is experiencing a huge impulse currently. And more specifically, Open Source robotics is growing rapidly. 
%\\

One of the open-source projects that appeared with this new way of thinking was the low-cost 3D printers. These machines have changed the manner of investigating in robots, since they allow to design them easily and have a 3D reproduction in a few hours. 
\\

But they are used not only in robots, but also in medicine and many other fields. In medicine, they have been a revolution since they allow to create customized pieces in very few time. They have been used for prosthesis and implants for persons of various ages, even for babies. 
\\

Before, the prosthesis were very expensive and permitted only fixed movements and combinations. The adaptations to each individual were made in the final product itself, trying to make it as comfortable as possible for the wearer. Nowadays, the prosthesis are customized for each patient, reducing the bothers for them and increasing their usability. Also, they can be easily and cheaply adapted for children as an example, who are still experimenting many changes in their bodies. They are much cheaper than they were before, and everyone with a 3D printer may construct one. 
\\

Of course, in order to 3D print a piece a file with its description is needed. There are many web-pages that store open-source designs that ranges from decoration models to complex prosthesis. This fact is decisive because there is not needed a huge amount of knowledge or money to improve the life quality of a person using these technologies. 
\\

There are numerous open-source projects and developers that put in common their knowledge to improve the technology being used. I have used many of them in the previous years, to learn about 3D printing, robotics or programming among other fields. 
\\

The fact that previously I have developed other projects involving designing and constructing 3D-printed wheeled robots gave me the idea of developing a new package for different robots that allowed to easily obtain new functionalities.One of the most interesting fields that can be used in virtually all types of robots is Computer Vision. 
\\

Computer vision has experience an important improvement in the last years through the upgrade of the hardware and software that compose it. The hardware such as acquisition elements (cameras, depth sensors, etc) and computing elements (PCs or other programmable devices) have experienced a rapid advance in the past years. It allowed to process more data that is now obtained more accurately and with less noise. This increase in the computing power of the equipment created a possibility of introducing more complex libraries and frameworks and even operating systems. 
\\

But Computer Vision is an ample area, why choosing object recognition? The project being developed is intended to be useful, not just a mere hobby. It occurred to me that apart from being useful in the interaction with robots and cognitive robots more particularly, this software could be used with visually impaired people, as an example. Mainly, with people that had a sudden loss of vision and that is not used to detect the objects by its shape, or cannot differentiate them by the contour. 
\\

In summary, the software developed is open-source and is intended as a useful package for visually impaired persons, as well as for its integration within a robotic system. 

\newpage
\addcontentsline{toc}{chapter}{Introduction}
\chapter*{Introduction}

% Intro

This bachelor's thesis consists on a software that implements an in-hand object learning and recognition using 2D and 3D information. 
\\

The code was designed to be running inside a robot. Hence, everything is modularized and all processes are designed to run in parallel. 

The hardware needed for this project is an RGB-D sensor and a computer. 

\addcontentsline{toc}{section}{What does the software do?}
\section*{What does the software do?}
% What does it do?
The input of the software is the image and the point cloud provided by the RGB-D sensor. The output of the system is different depending on the mode we are working. If the user is making it learn a new object, the output is the percentage of learning completed. On the other hand, if the user is utilizing the recognizing mode, the output of the system is the identifier of the object, i.e. the number that represents which object in the dataset seems to be more similar to the one being shown to it currently. 

%%%%%%%%%%%%%%%%%%%%%%%%%%%%%%%%%%PICTURE

\addcontentsline{toc}{section}{How can it be used?}
\section*{How can it be used?}
% How it is used?
From the beginning, a easy and intuitive human-machine interface was a must. In order to achieve this, this simple state 
machine was designed:

%%%%%%%%%%%%%%%%%%%%%%%%%%%%%%%%%%%%%PICTURE

The software only allows one hand being used at a time. It will select the hand that is located highest. 

In order to switch between the learning and recognizing modes of the program, the distance between the hand and the body
is being used: 

If the hand is stretched out towards the RGB-D sensor, the software will start learning. If it is closer to the body, 
it will launch the recognizing mode, which is the default. 

\addcontentsline{toc}{section}{General characteristics of the software}
\section*{General characteristics of the software}
% How is the code coded?

The Robotic Operating System is being used in this project since it has numerous features that allow a more efficient and 
easier implementation of software that has the characteristics described above. More information about this operating 
system may be found in the section.
%%%%%%%%%%%%%%%%%%%%%%%%%%(!!!!SECTION).

The software developed in this thesis is a ROS package. The last release might be downloaded from the following link. 
%%%%%%%%%%%%%%%%%%%%%%%%%%%%%%%%LINK


\addcontentsline{toc}{section}{Thesis structure}
\section*{ Thesis structure}
 
\newpage
%\addcontentsline{toc}{chapter}{Thesis structure}
\chapter{ Thesis structure}

This thesis is structured as follows. First, the state of the art of the different elements both hardware and software related with the project are presented. Afterwards, in the Methods part the project itself is detailed. The different key parts are thoroughly described and the reasons behind each selection of an element are explained. Then, the experiments made in the systems and its results are shown and they are compared and analysed in the following part, discussion. 
\\
Finally, the conclusions extracted from those experiments and the possible future lines of work are exposed in the last part of the project. 
After this last part, there appear the bibliography and references section, in which the --- used are cited, and the appendices. In the appendices it is possible to find additional documentation regarding the project such as a Gantt diagram or the minutes of meeting, and also the software requirements document or the software detailed documentation generated. 

